\section{Описание}

В данной лабораторной работе исследовалось качество реализации словаря на основе дерева PATRICIA, выполненной в предыдущей лабораторной работе. Основное внимание уделялось двум аспектам:
\begin{itemize}
    \item времени выполнения основных операций;
    \item корректности и объёму использования динамической памяти.
\end{itemize}

Для исследования были подготовлены отдельные профилировочные сценарии для операций \texttt{insert}, \texttt{search}, \texttt{remove}, \texttt{save} и \texttt{load}. Это позволило отделить собственную стоимость операций словаря от затрат на генерацию данных, сравнение с другими контейнерами и прочую обвязку тестовой программы.

\subsection{Использованные средства}

В качестве инструментов анализа были выбраны:
\begin{itemize}
    \item \texttt{gprof} для поиска наиболее дорогих по времени участков программы;
    \item \texttt{valgrind --tool=memcheck} для проверки корректности работы с динамической памятью;
    \item \texttt{valgrind --tool=massif} для оценки потребления памяти.
\end{itemize}

Таким образом, требование лабораторной работы о применении средств анализа было выполнено с использованием \texttt{gprof} и \texttt{Valgrind} как более универсальной замены \texttt{dmalloc}.

\subsection{Дневник выполнения}

\begin{longtable}{p{0.16\textwidth} p{0.78\textwidth}}
\toprule
Этап & Выполненные действия \\
\midrule
1 &
За базовую версию была взята реализация словаря из лабораторной работы 2. Для исследовательской части была создана отдельная директория лабораторной работы 3, чтобы не смешивать исходное решение и профилировочные материалы. \\

2 &
Был подготовлен бенчмарк для измерения времени выполнения операций словаря и созданы отдельные программы для профилирования операций \texttt{insert}, \texttt{search}, \texttt{remove}, \texttt{save} и \texttt{load}. \\

3 &
На наборе из \(200000\) слов были получены первые раздельные замеры времени: \texttt{insert} --- \(168886\) мкс, \texttt{search} --- \(131584\) мкс, \texttt{remove} --- \(113466\) мкс, \texttt{save} --- \(210991\) мкс. Для \texttt{load} на этом этапе был получен только предварительный результат, так как позднее методика измерения этой операции была уточнена отдельно. \\

4 &
При помощи \texttt{gprof} были сняты профили по операциям. Для \texttt{insert}, \texttt{search} и \texttt{remove} профиль оказался недостаточно детальным, так как при сборке с \texttt{-O2 -pg} существенная часть времени агрегировалась в \texttt{main}. Для \texttt{save} профиль оказался информативным и показал заметные затраты на работу временной хеш-таблицы и перераспределение памяти. \\

5 &
При помощи \texttt{valgrind --tool=memcheck} была проверена корректность работы с памятью. Ошибок доступа к памяти и утечек на исследуемых сценариях обнаружено не было. \\

6 &
При помощи \texttt{valgrind --tool=massif} была получена оценка пикового потребления памяти. Для широкого тестового сценария пик составил примерно \(16.95\) МБ, однако было отмечено, что заметная часть этой памяти относится к самой тестовой обвязке, а не только к словарю. \\

7 &
По результатам первого цикла профилирования была выполнена оптимизация операции \texttt{save}: из неё было убрано построение временного \texttt{unordered\_map<Node*, uint32\_t>}, а идентификаторы узлов стали временно храниться непосредственно в самих узлах. Одновременно был добавлен предварительный \texttt{reserve} для массива узлов. \\

8 &
После первой оптимизации время \texttt{save} уменьшилось с \(210991\) мкс до \(116424\) мкс на том же наборе из \(200000\) слов. Профиль \texttt{gprof} подтвердил исчезновение \texttt{unordered\_map::operator[]} и \texttt{rehash} из числа основных горячих точек. \\

9 &
Затем была улучшена методика измерения \texttt{load}: подготовка бинарного файла была вынесена из профилируемого процесса, что позволило получить более чистый профиль именно операции загрузки. \\

10 &
Дополнительный анализ показал, что значительная часть времени в \texttt{load} тратится на финальную проверку через повторный поиск каждого сохранённого ключа. В качестве второй оптимизации эта проверка была заменена на линейную структурную проверку загруженного дерева. Одновременно словарь начал хранить текущее число узлов, что позволило убрать ещё один лишний обход при сохранении. \\

11 &
После второй оптимизации были получены новые результаты: \texttt{save} --- \(75266\) мкс, \texttt{load} --- \(65724\) мкс для словаря из \(200000\) элементов. \\
\bottomrule
\end{longtable}

\subsection{Начальные времена основных операций}

Для наглядности в таблице ниже приведены первые раздельные замеры времени для основных операций словаря на наборе из \(200000\) слов.

\begin{center}
\begin{tabular}{lr}
\toprule
Операция & Время, мкс \\
\midrule
\texttt{insert} & \(168886\) \\
\texttt{search} & \(131584\) \\
\texttt{remove} & \(113466\) \\
\texttt{save} & \(210991\) \\
\bottomrule
\end{tabular}
\end{center}

Из этой таблицы видно, что на исходной версии наиболее дорогой операцией была сериализация словаря. Для \texttt{load} первоначальный замер не использовался в итоговом сравнении, так как после очистки методики для этой операции было получено более корректное базовое значение \(150250\) мкс. После этого стало видно, что и загрузка содержит заметные накладные расходы в завершающей проверке структуры. Именно поэтому оптимизация была сосредоточена прежде всего на операциях \texttt{save} и \texttt{load}.

\section{Найденные недочёты}

В ходе исследования были выявлены следующие недочёты:
\begin{itemize}
    \item исходный общий бенчмарк был недостаточно удобен для \texttt{gprof}, так как заметная часть времени приходилась на код вне самой структуры данных;
    \item в операции \texttt{save} имелись лишние накладные расходы из-за построения временной хеш-таблицы для отображения указателей в индексы и из-за повторных перераспределений памяти;
    \item в операции \texttt{load} применялась слишком дорогая финальная проверка корректности через повторный поиск каждого ключа;
    \item первоначальная методика профилирования \texttt{load} была не вполне чистой, потому что в том же процессе дополнительно подготавливался файл для загрузки;
    \item для \texttt{insert}, \texttt{search} и \texttt{remove} на использованной конфигурации профилирования не удалось выделить столь же выраженные внутренние узкие места, как для \texttt{save} и \texttt{load}, поэтому оптимизация была сосредоточена именно на сериализации и загрузке.
\end{itemize}

При этом операции \texttt{insert}, \texttt{search} и \texttt{remove} также были исследованы, однако для них не удалось получить настолько же локализованный и однозначно интерпретируемый результат профилирования, как для \texttt{save} и \texttt{load}.

\section{Сравнение версий}

В таблице приведены результаты для набора из \(200000\) слов.

\begin{center}
\begin{tabular}{lrrr}
\toprule
Операция & Исходная версия & После 1-й оптимизации & После 2-й оптимизации \\
\midrule
\texttt{save} & \(210991\) мкс & \(116424\) мкс & \(75266\) мкс \\
\texttt{load} & \(150250\) мкс & \(150250\) мкс & \(65724\) мкс \\
\bottomrule
\end{tabular}
\end{center}

Обе оптимизации затрагивали только внутренние накладные расходы реализации и не изменяли внешний формат словаря, набор поддерживаемых команд и их семантику.

Таким образом, обе выполненные оптимизации дали измеримый эффект. Первая оптимизация заметно ускорила сериализацию дерева, а вторая сократила время загрузки за счёт более дешёвой проверки структуры.

С точки зрения работы с памятью исправленная версия сохранила корректное поведение: \texttt{memcheck} не выявил ошибок и утечек, а измерения \texttt{massif} не показали ухудшения по памяти на исследованных сценариях. Следовательно, оптимизация проводилась не ради исправления ошибок доступа к памяти, а ради уменьшения внутренних накладных расходов.

\pagebreak
